\section{2024-10-22}

11-Meter ist kein Zufallsexperiment

\startitemize[packed]
\item Das Ergebnis darf nicht beeinflusbar sein
\item Es muss reproduzierbar sein
\stopitemize

11-Meter ist nur ein Zufallsexperiment, wenn es stark {\bf vereinfacht} wird.

\startframedtipp
Ein Versuch, bei dem alle Ergebnisse gleich Wahrscheinlich sind
\stopframedtipp

\startframedtipp
Wenn ein Zufallsexperiment mit einem Binärbaum darstellbar ist, ist es ein {\bf Bernoulli-Versuch}

\startitemize[packed]
\item Eigentlich einstufig
\item Kann wiederholt werden: in jeder Stufe muss der Selbe Versuch dargestellt werden. Das ist dan eine Bernoulli-Kette. Wenn man den Versuch $n$-Mal wiederholt ist es eine Kette von der länge $n$
\stopitemize
\stopframedtipp


\startbuffer[mp:bernoulli:kette:karten]
\startMPpage

numeric x, y, a;

% a = 2;
%
% label("(8 / 2**" & decimal a & ") = " & decimal (8 / (2**a)), origin);
% label(decimal 2*2, origin);


draw fullcircle scaled 20 shifted origin;
label(decimal 0 & decimal 0, origin);

for a=1 upto 3:
	for i=1 upto 2^a:
		x := ((240 / (2**a)) * (i - 0.5)) - 120;  % Horizontal spacing depends on level
		y := -40 * a;
		draw fullcircle scaled 20 shifted (x,y);
		label(decimal a & decimal i, (x,y));
	endfor
endfor

\stopMPpage
\stopbuffer

\placefigure[force][mp:bernoulli:kette:karten]{Darstellung einer bernoulli-Kette der Länge 3 für beim ziehen von Karten, ob es Kreuz ist, oder nicht.}{
\typesetbuffer[mp:bernoulli:kette:karten]
}

\startframedtipp
Klausuraufgaben kann man x y als bernoulli modelieren (S. 272 Nr. 6)
\stopframedtipp
